\section{The question left by ten percent}
\label{sec:question}
An averaging argument in the source selects a retained set of arithmetic positions. It can certify several possible sets on the same raw host. Earlier experiments asked whether a property of the selected witness followed from that selection argument. This led to a concrete way of constructing alternatives. Fix the host and its labels, specify how many points to draw from each of a small number of residue classes, and show that some draw passes the original selector's score. Once selected, use that one draw across all the primes, lifts and comparisons that define the gain. If one independently picks a different retainer after seeing each prime, local recovery can look large even though there is no common arithmetic witness.

The present question concerns a second order of choices. Before reading the labels, is there a \emph{rule} guaranteed to recover ten percent of the odd-position allowance? The previous study gave such a rule and a hand certificate with very little slack. What was tight in that calculation? At least three explanations deserve a trial. The assigned rates may have put too many positions in expensive residues. The fixed support may have allowed too many adjacent pairs. Or the numerical order itself may have left every geometrically feasible support with expensive odd positions. Each explanation predicts a different response to intervention. Moving mass within one fixed support tests the first, replacing supports at fixed resolution tests the second, and constructing an adverse order \emph{against every permitted replacement} tests the third.

\subsection{Three optimization problems and their witnesses}
A labeling induces a preference tree \(L\). In each cell the child containing its least host label is called its preferred or canonical even child. A profile \(a=(a_t)_{t=0}^{63}\) specifies the asymptotic rate in each additive residue \(t\), measured in units of the source density \(\delta\). The exact-count normalization is \(\sum a_t=64\). Let \(S=\{t:a_t>0\}\). Later sections give an integer-coefficient numerator \(P_L(S,a)\) and a strictly positive denominator \(Q_L(a)\). Write \(f_L=P_L/Q_L\). Here \(P\) credits odd mass at a cell for each missing generator class. The quantity \(Q\) charges odd mass for all generator classes, with a within-residue tail. Both use the same rates.

Let \(\mathcal A(L,S)\) be the profiles passing a specified sufficient tile test on \(S\). A policy \(S_0(L)\) that first names the occupied leaves defines the restricted value
\[
 V_0=\min_L\ \sup_{a\in\mathcal A(L,S_0(L))} f_L(S_0(L),a).
\]
One can also allow the policy to choose its support:
\[
 V_{\rm tile}=\min_L\ \sup_{S}\ \sup_{a\in\mathcal A(L,S)} f_L(S,a).
\]
There are finitely many trees at this resolution. A strict score margin gives an open rate condition, hence the supremum. Specifying an interior rational profile proves a lower bound. A limiting upper value can remain unattained at any fixed host size or any one strict rational certificate.

For a third problem replace the quantitative mark test by a cap on individual rates and the necessary condition that an occupied node has an occupied preferred child. This gives an outer model \(V_{\rm box}(c)\). A larger admissible set raises the inner optimum, so a single labeling with box optimum \(u\) proves \(V_{\rm tile}\le u\) whenever the tile class lies within that box. A construction on just one labeling supplies a lower bound on that labeling's own maximum. It supplies no worst-label lower bound. Conversely, an arbitrary feasible construction on an adverse labeling supplies no worst-label upper bound. An upper bound requires control over \emph{all} admissible replies to the adversary.

Here every labeling means every abstract numerical order of the residue tree. A tree can always be realized by listing the preferred subtree first recursively. Whether a specified adverse order occurs along the underlying number-field construction is a different inverse problem. A lower guarantee proved on all trees applies immediately to any smaller family of actual trees. A counterexample on one abstract tree only rules out a larger label-independent theorem in the specified model. This is why we publish the adverse word itself.

\subsection{The unit of the percentage}
If the allowance has normalized weight \(K\) and the fallback recovers \(R\), the statistic is \(R/K\). A claim of ten percent does not add ten percentage points to the slope gap. An absolute change in that gap additionally depends on the size of \(K\) and the other ledger entries. The gain and the charge are compared at the same prime and lift before averaging. Replacing the denominator by the universal half-allowance would answer another question and can reorder two policies.

This distinction played a role earlier. A bound on the weight of odd positions limited the loss that could enter a column. The gain identity then showed that a part of this potential loss is recovered. A lower construction showing positive recovery therefore defeats the proposition that selection alone makes recovery negligible. Sharpening \(V_{\rm tile}\) refines the size of that obstruction. It neither identifies the set returned by classical choice nor locates a zeta zero. The two witnesses, the original existential one and the constructed alternative, play different roles in the inference.

\subsection{Tests that discriminate the explanations}
The first test keeps the old decision on every labeling and recomputes its rate allocation. An upper certificate on the most difficult decision makes that intervention interpretable. A second test changes the actual retained leaves and covers the resulting relative-label classes. It records how many classes clear each target, including failures. Coverage is a combinatorial count. A probabilistic interpretation would require a law on arithmetic labelings and its multiplicities under normalization.

The third test gives the adversary the first move and lets the responding policy use every support in a larger finite class. The search for a difficult word may be heuristic. Once the word has been written down, the proof that it is difficult must be exhaustive. Finally we weaken the calculation by discarding additive positions and retaining only subtree counts. This last experiment tests why the earlier calculations work. A large jump in the relaxed ceiling measures information lost by that description. It does not measure an improvement available to the selector.
